\section{Conclusion}
\label{sec:clustering_agent_representation_conclusion}

The \gls{car} method was developed to investigate how concepts are organized within the representations of \gls{rl} policies.
It is based on three hypotheses introduced throughout this manuscript, the \MechanisticHypothesisName, the \ConceptHypothesisName, and the \RepresentationConvergenceHypothesisName.

\gls{car} translates these assumptions into an empirical analysis.
Multiple surrogate policies are trained independently to reproduce the same policy.
Their \gls{lr} are then clustered and compared to assess whether similar conceptual structures emerge across policies.
The \ClusteringUniversalityName and \AbstractionScoreName provide quantitative measures for evaluating the consistency and abstraction of the resulting clusters.

This makes \gls{car} a general procedure for analyzing agent representations.
The analysis and evaluation metrics are fully automated and require no manual intervention or environment-specific adjustments.
The method can therefore be applied to a wide range of \gls{rl} settings and provides a basis for systematically studying the internal representations of neural policies.

The next chapter applies \gls{car} to several environments to evaluate its ability to identify meaningful conceptual structures.

% A more fundamental limitation concerns the semantic interpretation of the clusters.
% While \gls{car} provides a quantitative description of the structure present in the agent's representations, the meaning assigned to individual clusters still relies on manual interpretation.
% \gls{car} therefore characterizes how the \gls{ls} is organized, but not what each concept stands for.
% Developing methods to automatically extract the semantics of each cluster is an important direction for future work.

% When situating \gls{car} within the field of explainability, substantial work remains to be done.
% The structure it recovers can be considered an explanation in the sense defined in Section~\ref{sec:explanation}, since it provides information about the concepts represented in the agent's \gls{ls}.
% How such explanations can be used to support validation, diagnosis, or certification, as discussed in Chapter~\ref{chap:explainability_reinforcement_learning}, remains to be established.
% Connecting \gls{car} explanations more directly to the agent's decisions and failures would help determine how they can serve these broader objectives.
